\documentclass[11pt]{article}
\usepackage[margin=1in]{geometry}
\usepackage{amsmath,amssymb,amsthm}
\usepackage{enumitem}
\usepackage{hyperref}
\usepackage{booktabs}

\newcommand{\code}[1]{\texttt{#1}}

\title{Guide to Interpreting the Output of the\\
Segmented Fourier Analysis Code (Version~9.4)}
\author{}
\date{}

\begin{document}
\maketitle

\section{Overview}

The v.\,9.4 code takes as input a list of pairs $(n, \mu_n)$, where
$\mu_n$ is the minimum modulus among the eigenvalues of a certain matrix
associated to the index~$n$.  Its purpose is to determine whether the
sequence $\{\mu_n\}$ oscillates in a nearly periodic fashion and, if so,
to estimate the period.

\subsection*{What changed from earlier versions}

Versions~8 and 9 attempted to remove any overall trend (growth, decay,
curvature) from the data by fitting a single curve to the entire dataset
before analyzing the residual for periodicity.  This failed on data
composed of qualitatively different regions---for instance, a linearly
rising segment followed by a flat plateau.  No single polynomial or
smoother can follow the sharp bend between such regions, and the resulting
artifact corrupts all downstream analysis.

Version~9.4 takes a fundamentally different approach:
\begin{enumerate}
\item \textbf{Segmentation.}  Automatically detect \emph{breakpoints}
  where the data's behavior changes, and split the data into segments at
  those points.
\item \textbf{Per-segment analysis.}  Within each segment, the data is
  approximately linear, so a simple straight-line fit suffices to remove
  the trend.  Each segment is then analyzed independently for periodicity.
\end{enumerate}

The report therefore contains one verdict per segment rather than a single
verdict for the whole dataset.

\section{Glossary of abbreviations and acronyms}
\label{sec:glossary}

\begin{description}[style=nextline, leftmargin=3cm]

\item[ACF] \textbf{Autocorrelation function.}  Given a sequence
  $x_1, \ldots, x_n$ with mean $\bar{x}$, the normalized autocorrelation
  at lag~$k$ is
  \[
    \mathrm{ACF}(k)
      = \frac{\sum_{i=1}^{n-k}(x_i - \bar{x})(x_{i+k} - \bar{x})}
             {\sum_{i=1}^{n}(x_i - \bar{x})^2}.
  \]
  By construction $\mathrm{ACF}(0) = 1$.  If the sequence repeats with
  period~$T$, the ACF has a local maximum near~$1$ at lag~$T$.
  See any introductory time series text, or the Wikipedia article
  \emph{Autocorrelation}, for further background.

\item[ACF strength] The value of $\mathrm{ACF}(T)$ at the detected
  period~$T$.  Ranges from $-1$ to $+1$.  Values near~$+1$ indicate
  strong, clean repetition.  The code uses $0.5$ as the threshold for a
  confident detection and $0.2$ as the threshold below which
  the detection is dismissed.

\item[AR(1)] \textbf{First-order autoregressive model.}  A stochastic
  process in which each value depends linearly on its predecessor:
  \[
    x_n = \rho\, x_{n-1} + \varepsilon_n, \qquad
    \varepsilon_n \sim \mathcal{N}(0, \sigma^2).
  \]
  The parameter $\rho$ (``rho'') is the lag-1 autocorrelation.
  This model serves as the \emph{null hypothesis}: it represents data that
  has short-range correlation but no true periodicity.  A spectral peak
  is declared significant only if it exceeds what the AR(1) model would
  produce.  See Gilman, Fuglister, and Mitchell, ``On the power spectrum
  of `red noise'\,'' \emph{J.\ Atmos.\ Sci.}\ \textbf{20} (1963),
  182--184, for the theoretical power spectrum used.

\item[Bonferroni correction] A multiple-testing correction.  When
  $m$~frequency bins are tested simultaneously at significance
  level~$\alpha$, up to $m\alpha$ false positives are expected by chance.
  The Bonferroni correction tests each bin at level $\alpha/m$, so that
  the probability of \emph{any} false positive across all bins stays
  at~$\alpha$.  Conservative but simple and widely accepted.

\item[Breakpoint] An index $n$ at which the data's local behavior changes
  qualitatively---for example, from rising to flat, or from one slope to
  another.  Detected by the PELT algorithm (see below).

\item[dB] \textbf{Decibel.}  $10\log_{10}(\text{power})$.  A logarithmic
  unit for power, used in some plots.

\item[Dynamic range] The ratio of the 90th percentile to the 10th
  percentile of the estimated envelope.  If this exceeds~3, the code
  concludes that the oscillation amplitude varies enough to require
  normalization.

\item[Envelope] The slowly varying curve tracing the peaks of an
  oscillating signal.  For $A(n)\sin(2\pi n/T)$, the envelope
  is~$|A(n)|$.

\item[FFT] \textbf{Fast Fourier Transform.}  An $O(n\log n)$ algorithm
  for computing the discrete Fourier transform of a length-$n$ sequence,
  decomposing it into sinusoidal components at frequencies $k/n$,
  $k = 0, 1, \ldots, n-1$.  See any numerical analysis text or the
  Wikipedia article \emph{Fast Fourier transform}.

\item[Fundamental period] The true repeat period of the data, as
  distinguished from harmonics (integer fractions of the period) and
  integer multiples (longer lags at which the ACF also peaks).  Identified
  by the multi-scale ACF procedure described in
  Section~\ref{sec:multiscale}.

\item[Hann window] A bell-shaped taper $w_k = \frac{1}{2}
  \bigl(1 - \cos(2\pi k / (n-1))\bigr)$ applied to the data before the
  FFT\@.  It reduces \emph{spectral leakage} (spurious spreading of power
  across frequencies caused by the finite length of the data) but
  attenuates signals near the Nyquist frequency.

\item[Harmonic] An integer multiple of a fundamental frequency.  If the
  fundamental frequency is $f_0 = 1/T$, the $m$-th harmonic is $mf_0$,
  corresponding to period $T/m$.  Non-sinusoidal waveforms (e.g.,
  sawtooth, square wave) have energy at many harmonics; the FFT reports
  these as separate peaks.

\item[MA($w$)] \textbf{Moving average with window width~$w$.}  The value
  at index~$i$ is replaced by the mean of the $w$ values centered at~$i$.
  Used in the multi-scale period detection to suppress fine-scale structure.

\item[Nyquist frequency] The highest frequency representable by a
  discrete sequence: $f_{\mathrm{Nyq}} = 1/2$ (cycles per sample),
  corresponding to a period of~2.  The Hann window attenuates signals
  near this limit, which is why a period-2 oscillation can be missed by
  the FFT while the ACF detects it.

\item[$n/3$] The maximum period considered credible within a segment of
  length~$n$.  A period of $T$ in $n$~points gives $n/T$ full cycles;
  for $T > n/3$ there are fewer than three full cycles, which is
  insufficient to distinguish periodicity from a trend artifact.

\item[PELT] \textbf{Pruned Exact Linear Time.}  A changepoint detection
  algorithm that finds the positions and number of breakpoints in a
  sequence by minimizing a penalized cost function.  The penalty
  (proportional to $\log(N) \cdot \operatorname{Var}(\mu)$ in this code)
  controls sensitivity: a larger penalty produces fewer breakpoints.
  Reference: Killick, Fearnhead, and Eckley, ``Optimal detection of
  changepoints with a linear computational cost,''
  \emph{J.\ Amer.\ Statist.\ Assoc.}\ \textbf{107} (2012), 1590--1598.
  The implementation used is from the \code{ruptures} Python library
  (Truong, Oudre, and Vayatis, 2020).

\item[Power] The squared magnitude $|c_k|^2$ of a Fourier coefficient,
  measuring how much of the signal's variance is attributable to
  oscillation at frequency~$k/n$.  Displayed in scientific notation:
  \code{1.58e+04} means $1.58 \times 10^4$.

\item[Power spectrum] The function mapping each frequency to its power.
  A peak in the power spectrum indicates energy concentrated at that
  frequency.

\item[rho ($\rho$)] The lag-1 autocorrelation of the analysis signal,
  used as the AR(1) parameter.  $\rho$ near $+1$: strong positive
  correlation.  $\rho$ near $-1$: alternation.  $\rho$ near $0$:
  near-independence.

\item[Savitzky--Golay filter] A smoothing method that fits a low-degree
  polynomial by least squares within a sliding window, replacing each
  value with the fitted value.  Used here for envelope estimation.
  Reference: Savitzky and Golay, \emph{Anal.\ Chem.}\ \textbf{36}
  (1964), 1627--1639.

\item[Spectral leakage] Spurious spreading of a signal's power across
  neighboring frequencies, caused by the FFT's implicit assumption that
  the data is periodic over its length.  Reduced by applying a Hann
  window before the FFT.

\item[std dev] Standard deviation.
\end{description}

\section{Section~1 of the report: data summary}

The report begins with the range of the index~$n$, the number of data
points~$N$, and the range of the minimum moduli $\mu_n$.

\section{Section~2 of the report: segmentation}
\label{sec:segmentation}

The code approximates the data by a piecewise-linear function---that is,
a sequence of straight line segments joined at breakpoints---and finds the
breakpoint positions that minimize the total squared error, subject to a
penalty for each additional breakpoint.  This is done by the PELT
algorithm (see glossary).

The report lists:
\begin{itemize}
\item the number of breakpoints detected,
\item for each resulting segment: the index range, the number of points,
  and the range of $\mu_n$ values within that segment.
\end{itemize}

\noindent\textbf{Example.}  Data that rises linearly from $n=1$ to
$n=120$ and then levels off would produce one breakpoint near $n=120$,
giving two segments: the ramp (Segment~1) and the plateau (Segment~2).

If no breakpoints are detected, the entire dataset is treated as a single
segment.

\section{Section~3 of the report: per-segment analysis}

Each segment is analyzed independently.  The analysis has five stages,
described in the subsections below.  Segments shorter than 10~points are
skipped.

\subsection{Linear detrend}

A straight line is fitted to the segment by least squares and subtracted.
Within a single segment, the data is approximately linear (that is the
whole point of segmenting), so this simple detrend is adequate.  The
report prints the slope and the standard deviation of the residual.

\subsection{Envelope normalization}

If the oscillation amplitude varies strongly within a segment (for
instance, an oscillation that decays as $1/n$), the ACF and significance
tests become unreliable.  The code estimates the local amplitude envelope
by smoothing the absolute values of the detrended signal with a
Savitzky--Golay filter, then checks the \emph{dynamic range} (90th
percentile divided by 10th percentile of the envelope).  If the dynamic
range exceeds~3, the code divides the detrended signal by the envelope,
producing a signal with roughly constant amplitude.

When normalization is applied, the report notes this together with the
dynamic range.  When it is not needed, the report says so.  All subsequent
analysis operates on the signal after this step, called the
\emph{analysis signal}.

\subsection{Windowed FFT}

The analysis signal is multiplied by a Hann window and passed through the
FFT, yielding a power spectrum.  See the glossary entries for FFT, Hann
window, power spectrum, and spectral leakage.

\subsection{Significance testing}
\label{sec:significance}

To determine whether a peak in the power spectrum represents genuine
periodicity rather than noise, the code constructs an AR(1) null model
(see glossary) fitted to the analysis signal.  The theoretical power
spectrum of this null model, which is not flat but tilted toward low or
high frequencies depending on the sign of~$\rho$, serves as the baseline.

The 95\% and 99\% confidence thresholds are computed from a chi-squared
distribution with 2~degrees of freedom (the approximate distribution of
FFT power under the null), then inflated by the Bonferroni correction
(see glossary) to account for the number of frequency bins tested.

Peaks exceeding the threshold are listed in a table with their period,
power, and significance level (95\% or 99\%).

\subsection{Multi-scale ACF period detection}
\label{sec:multiscale}

The FFT can misidentify the fundamental period of non-sinusoidal
waveforms, because sharp features (sawtooth teeth, square-wave edges)
distribute energy across many harmonics, and a harmonic can exceed the
fundamental in power.  The multi-scale ACF procedure avoids this by
working directly with the autocorrelation, which is insensitive to
waveform shape.

The procedure examines the analysis signal at progressively coarser
scales by smoothing it with moving averages of increasing width.  At each
scale, it computes the ACF and records the first peak whose strength
exceeds~$0.1$.  Fine scales detect short-period features; coarse scales
detect longer periods after the short-period structure has been averaged
out.

\subsubsection*{Identifying the fundamental}

Among all detected periods, the code selects the one whose integer
multiples best account for the others.  For each candidate period~$T_c$,
it counts how many of the other detected periods are approximately equal
to $mT_c$ for some positive integer~$m$ (within a 15\% tolerance).  The
candidate explaining the most others wins; ties are broken by ACF
strength.

This correctly handles two situations:
\begin{itemize}
\item A period-2 oscillation whose ACF peaks at lags $2, 4, 6, 8$:
  period~2 explains all the others as multiples, so it is the fundamental.
\item A period-60 sawtooth with sub-cycle teeth of period~3: period~60
  and period~3 may both appear, but the tie is broken by ACF strength,
  which favors the coarser-scale detection of period~60.
\end{itemize}

\section{Verdicts}

Each segment receives one of four verdicts:

\begin{description}[style=nextline, leftmargin=3cm]
\item[PERIODIC] At least one of the following holds: (a)~the FFT has a
  peak exceeding the Bonferroni-corrected 95\% or 99\% threshold, or
  (b)~the multi-scale ACF finds a fundamental with strength above~$0.5$.
  The fundamental period and ACF strength are reported.

\item[PERIODIC (FFT only)] The FFT has significant peaks, but the
  multi-scale ACF does not find a strong fundamental.  The dominant FFT
  period is reported.

\item[POSSIBLY PERIODIC] The multi-scale ACF finds a candidate period
  with strength between $0.2$ and~$0.5$, but the FFT does not confirm.
  The candidate period is reported with a caution.

\item[NOT PERIODIC] Neither the FFT nor the ACF finds evidence of
  oscillatory behavior in this segment.
\end{description}

\section{Overall summary}

After all segments have been analyzed, the report prints a one-line
summary for each segment: the index range, the verdict, and (if
applicable) the fundamental period and ACF strength.  This is the primary
output for quick reference.

\section{Visualizations}

\subsection{Overview plots (one figure)}

\begin{enumerate}
\item \textbf{Original Data with Detected Breakpoints.}  The full dataset
  with vertical red dashed lines at each breakpoint.  Verify that the
  breakpoints fall where the data's character visibly changes.

\item \textbf{Segment-wise Detrended Residuals.}  The residuals from each
  segment's linear detrend, plotted in different colors.  Each segment's
  residual should oscillate around zero with no visible drift.  If a
  segment's residual has a large broad hump or slope, the segmentation may
  have missed a breakpoint.
\end{enumerate}

\subsection{Per-segment detail plots (one figure per segment)}

Each segment produces four panels:

\begin{enumerate}
\item \textbf{Data with Linear Trend.}  The segment's data (blue) with the
  fitted line (red).  The line should follow the broad slope of the data
  without tracking oscillations.

\item \textbf{Analysis Signal.}  The signal that the FFT and ACF operate
  on (after detrending and, if applicable, envelope normalization).  It
  should oscillate around zero with roughly constant amplitude.

\item \textbf{Power Spectrum.}  Log-scale plot of power versus frequency.
  The dashed lines are the AR(1) null model (black), the 95\%
  Bonferroni-corrected threshold (green), and the 99\% threshold (red).
  Colored dots mark peaks exceeding a threshold.

\item \textbf{Power vs Period.}  The same information as the power
  spectrum, but with the horizontal axis showing period (the reciprocal of
  frequency) rather than frequency.  Restricted to periods between~2 and
  $n/3$.
\end{enumerate}

\section{Practical notes}

\subsection*{When the segmentation is wrong}

The PELT algorithm's penalty parameter controls sensitivity.  The code
sets it to $\log(N) \cdot \operatorname{Var}(\mu)$, which is a standard
default.  If you believe a breakpoint was missed or spuriously inserted,
the penalty can be adjusted by editing the line
\code{penalty\_value = np.log(N) * np.var(minmoduli)} in the code.
Increasing the penalty produces fewer breakpoints; decreasing it produces
more.

\subsection*{When a segment has no breakpoints}

If the data is already well-described by a single straight line (or is
roughly flat with oscillations), no breakpoints will be detected, and the
entire dataset is analyzed as one segment.  The analysis within that
segment is identical to what earlier versions did (minus the complications
of global polynomial or Savitzky--Golay detrending).

\subsection*{Dependencies}

Version~9.4 requires the \code{ruptures} Python library, which is not
included in SageMath by default.  Install it with
\code{!pip install ruptures} in a SageMath notebook cell, or
\code{pip install ruptures} from the command line.

\end{document}
